% =====================================================================
%  c02：原印刷页 14–21（PDF p.15–22）
%  内容：§2.1.1 页末编号列表的第 2–4 项（第 1 项由 c01 末尾引出）、§2.1.2、
%        §2.1.3、§2.2、§2.2.1、§2.2.2（至 p.21 页末的 Remarks 标题）。
%  本片段从编号列表中部起，故不补标题；全片段没有 \section 级标题，
%  因此不插入 keytakeaways。本片段无插图；编号公式 (2.19)–(2.49) 连续。
% =====================================================================

\begin{enumerate}[start=2]
  \item 在量子理论中，当 $2\mathcal{B}=\theta/\pi$ 不是整数时，对称性破缺为 $SO(2)$。

  \item 在量子理论中，当 $2\mathcal{B}=\theta/\pi$ 为偶数时，理论仍然具有 $O(2)$ 对称性。

  \item 在量子理论中，当 $2\mathcal{B}=\theta/\pi$ 为奇数时，经典的 $O(2)$ 关系不再成立，但覆叠群 $\mathrm{Pin}^{+}(2)$ 是一个对称性。
\end{enumerate}

我们要强调，圆环上的粒子\underline{并不是}自旋二分之一！！话虽如此，如果我们定义角动量 $\mathcal{L}$，使得 $H=\frac{\mathcal{L}^{2}}{2I}$，那么当 $\mathcal{B}$ 为半整数时，角动量确实具有半整数本征值，一如人们对自旋表示所期待的那样。所以我们得到的是：半通量量子（half flux quantum）在诱导这个系统的半整数自旋，从而经典系统的经典 $O(2)$ 对称性在量子理论中被实现为 $\mathrm{Pin}^{+}(2)$ 对称性。这是出现在带非平凡规范场与 Chern-Simons 项的量子对称性中的一个耐人寻味的现象：粒子的统计与自旋可以相对其经典值发生移动，而且往往是以牵涉古怪分数的形式移动。

\subsubsection{围绕螺线管的圆环上的带电粒子：路径积分} % 原文: 2.1.2 Charged Particle On A Circle Surrounding A Solenoid: Path Integrals

用两种办法计算这个理论的欧氏时间传播子会很有启发性。于是我们考虑量 $\langle\phi_2|e^{-\beta H}|\phi_1\rangle$。一方面，我们已经把哈密顿量完全对角化，所以插入一组完备态便立刻得到：

\begin{equation}
  \langle\phi_2|e^{-\beta H_{\mathcal{B}}}|\phi_1\rangle=\frac{1}{2\pi}\sum_{m\in\mathbb{Z}}e^{-\frac{\beta}{2I}(m-\mathcal{B})^{2}+\mathrm{i}m(\phi_1-\phi_2)}
  \tag{2.19}
  \label{eq:2.19}
\end{equation}

% [Tutor Note: (2.19) 最容易出错的是相位里 $(\phi_1-\phi_2)$ 的次序与符号。用 $\Psi_m(\phi)=\frac{1}{\sqrt{2\pi}}e^{\mathrm{i}m\phi}$（见 \eqref{eq:2.11}）作完备性插入：$\langle\phi_2|e^{-\beta H_{\mathcal{B}}}|\phi_1\rangle=\sum_m\langle\phi_2|m\rangle e^{-\beta E_m}\langle m|\phi_1\rangle=\frac{1}{2\pi}\sum_m e^{\mathrm{i}m\phi_2}e^{-\beta E_m}e^{-\mathrm{i}m\phi_1}$，正是 $e^{\mathrm{i}m(\phi_1-\phi_2)}$。若把 $\langle\phi_2|\cdots|\phi_1\rangle$ 的左右次序写反，就会得到 $e^{\mathrm{i}m(\phi_2-\phi_1)}$，而 (2.19) 与 (2.23) 之间的一次 theta 变换核对将失败；这个符号也正是下一段 (2.20) 里作者提醒「必须小心虚项的符号，而且它很要紧」的那个符号。]

另一方面，利用传播子的 Feynman 路径积分表示作 Wick 转动，我们也可以得到一个用路径积分表达的式子：

\begin{equation}
  \langle\phi_2|e^{-\beta H}|\phi_1\rangle=Z(\phi_2,\phi_1|\beta):=\int[d\phi(t)]_{\phi(0)=\phi_1}^{\phi(\beta)=\phi_2}e^{-\frac{1}{\hbar}\int_0^{\beta}\frac{1}{2}I\dot\phi^{2}dt-\mathrm{i}\int\mathcal{B}\dot\phi dt}
  \tag{2.20}
  \label{eq:2.20}
\end{equation}

注意，在转到欧氏号差的这一步中，拓扑项仍然是虚的。必须小心取准这个虚项的符号，而它很要紧。

把这个理论看作场论，那么它是一个自由场论，其路径积分可以用半经典技巧精确计算：

运动方程简单地就是 $\ddot\phi=0$。同样，$\theta$ 项并没有改变它。

于是，带有边界条件 $\phi(0)=\phi_1$、$\phi(\beta)=\phi_2$ 的运动方程的经典解为：

\begin{equation}
  \phi_{cl}(t)=\phi_1+\left(\frac{\phi_2-\phi_1+2\pi w}{\beta}\right)t
  \qquad w\in\mathbb{Z}
  \tag{2.21}
  \label{eq:2.21}
\end{equation}

或者更切题地说：

\begin{equation}
  e^{\mathrm{i}\phi_{cl}(t)}=e^{\mathrm{i}\left(\left(1-\frac{t}{\beta}\right)\phi_1+\frac{t}{\beta}\phi_2\right)+\frac{2\pi\mathrm{i}tw}{\beta}}
  \tag{2.22}
  \label{eq:2.22}
\end{equation}

这些是欧氏运动方程的解，由于历史的原因被称为「瞬子」（instantons）。注意，由于我们定义这个 $0+1$ 维场论时所处的时空是紧致的，这样的解有无穷多个，以 $w\in\mathbb{Z}$ 标记。这个游戏里有两个圆：这个 $0+1$ 维场论的时空是欧氏时间圆；而场论的目标空间也是一个圆。量子的场 $e^{\mathrm{i}\phi(t)}$ 是一个映射 $M\to\mathcal{X}$。$M$ 是时空，它是 $S^{1}$；而 $\mathcal{X}$ 是目标空间，它也是 $\mathcal{X}$。回忆 $\pi_1(S^{1})\cong\mathbb{Z}$。可以存在拓扑上不等价的场构型。也就是说，映射空间 $Map(M\to\mathcal{X})$ 有不同的连通分支。不同的拓扑扇区由映射 \eqref{eq:2.22} 的绕数（winding number）唯一地标记。在路径积分里我们对所有场构型求和，因此也应当对所有这些瞬子构型求和。

% [Tutor Note: 原文此句「$\mathcal{X}$，也就是目标空间，也是 $\mathcal{X}$」显然是笔误，依据上一句应读作「也是 $S^{1}$」；照原文转写，未改。另需注意 $e^{\mathrm{i}\phi}$ 是映射 $M\to\mathcal{X}$，而不是 $\phi$ 本身——$\phi$ 只是 $\mathbb{R}$ 值的一个「抬升」，它在绕 $M$ 走一圈时只能回到模 $2\pi$ 的同一个值，绕数 $w$ 正是这个抬升的单值性障碍，也就是 $M$ 的 $\pi_1$ 到 $\mathcal{X}=S^{1}$ 的 $\pi_1$ 的映射度。]

一个直接的计算（细节见在线讲义 [48]）表明，路径积分表达式恰好等于

\begin{equation}
  Z(\phi_2,\phi_1|\beta)=\sqrt{\frac{I}{2\pi\beta}}\sum_{w\in\mathbb{Z}}e^{-\frac{2\pi^{2}I}{\beta}\left(w+\frac{\phi_2-\phi_1}{2\pi}\right)^{2}-2\pi\mathrm{i}\mathcal{B}\left(w+\frac{\phi_2-\phi_1}{2\pi}\right)}
  \tag{2.23}
  \label{eq:2.23}
\end{equation}

% [Tutor Note: 此处作者以「a straightforward computation（一个直接的计算）」跳过全部推导，只指向在线讲义 [48]。被略去的是两件事：一是围绕 (2.21)–(2.22) 的经典解作二次涨落积分，涨落算符为 $-\frac{I}{\hbar}\partial_t^{2}$（配以 Dirichlet 边界条件），其行列式给出前置因子 $\sqrt{I/2\pi\beta}$——(2.19) 中已取 $\hbar=1$，与 (2.20) 里显式的 $1/\hbar$ 对照即可定出这一因子；二是把「按 $w$ 求和」与「按 $m$ 求和」联系起来的 Poisson 求和公式，这才是 (2.23) 与 (2.19) 相等这一「看似神奇」的数学事实的真正内容，也正是紧接着 (2.24)–(2.25) 要给出的 theta 函数机制。核对符号时请记住：拓扑项给出的相位是 $-2\pi\mathrm{i}\mathcal{B}w$，其符号与 (2.20) 中的 $-\mathrm{i}\int\mathcal{B}\dot\phi dt$ 一致。]

现在比较 \eqref{eq:2.19} 与 \eqref{eq:2.23}。这两个表达式看上去大不相同！一个涉及以 $\beta$ 在分子上的指数之和，另一个则是 $\beta$ 在分母上。前者适合讨论 $\beta\to\infty$（低温）的渐近行为，后者适合讨论 $\beta\to0$（高温）。然而，我们算出的是同一个物理量，只是用了两种不同的方法，所以它们必定相等。但那个说它们相等的数学恒等式看起来颇有几分神奇。下面我们说明如何用一个直接的数学论证来验证这两个表达式确实完全相同。

我们引入 Riemann theta 函数

\begin{equation}
  \vartheta\left[\genfrac{}{}{0pt}{}{\theta}{\phi}\right](z|\tau):=\sum_{n\in\mathbb{Z}}e^{\mathrm{i}\pi\tau(n+\theta)^{2}+2\pi\mathrm{i}(n+\theta)(z+\phi)}
  \tag{2.24}
  \label{eq:2.24}
\end{equation}

Riemann theta 函数是 $\tau$ 在上半平面内的绝对收敛解析函数，它也是 $z$ 的整函数。利用 Poisson 求和公式可以验证关键的模变换律：

\begin{equation}
  \vartheta\left[\genfrac{}{}{0pt}{}{\theta}{\phi}\right]\left(\frac{-z}{\tau}\Big|\frac{-1}{\tau}\right)=(-\mathrm{i}\tau)^{1/2}e^{2\pi\mathrm{i}\theta\phi}e^{\mathrm{i}\pi z^{2}/\tau}\,\vartheta\left[\genfrac{}{}{0pt}{}{-\phi}{\theta}\right](z|\tau)
  \tag{2.25}
  \label{eq:2.25}
\end{equation}

% [Tutor Note: 「用 Poisson 求和公式可以验证」跳过的正是全部计算：把 $\sum_n\delta(x-n)$ 展成 $\sum_k e^{2\pi\mathrm{i}kx}$ 后再逐项作高斯积分，才会同时得到前置因子 $(-\mathrm{i}\tau)^{1/2}$、指数上的 $e^{\mathrm{i}\pi z^{2}/\tau}$，以及特征（characteristic）的交换 $\left[\genfrac{}{}{0pt}{}{\theta}{\phi}\right]\to\left[\genfrac{}{}{0pt}{}{-\phi}{\theta}\right]$（注意这不是简单的 $\theta\leftrightarrow\phi$，两处各带一个负号）。$(-\mathrm{i}\tau)^{1/2}$ 的平方根分支不可随意取：$\tau$ 在上半平面时须取主分支，否则 (2.25) 会整体多一个负号。用 (2.25) 核对传播子时，要点在于 $\tau$ 在此取纯虚值（$\tau\propto\mathrm{i}\beta$），于是模变换 $\tau\to-1/\tau$ 恰好把 $\beta\to0$ 与 $\beta\to\infty$ 互换——这正是 (2.19) 适合低温、而 (2.23) 适合高温的根源；同时还要把 $\mathcal{B}$ 与 $\frac{\phi_2-\phi_1}{2\pi}$ 分别对号入座地塞进特征 $(\theta,\phi)$，否则 (2.19) 的 $m$ 求和与 (2.23) 的 $w$ 求和很容易被互换，而丢掉 $2\pi\mathrm{i}\mathcal{B}(\cdots)$ 这一相位。]

可以很容易地核验：上面给出的欧氏传播子的两个公式，正是由 Riemann theta 函数的模变换联系起来的。

我们将会看到，theta 函数及其模性质在后文将起重要作用。

\medskip\noindent\textbf{练习}\quad 设 $2\mathcal{B}$ 为奇数。

\begin{enumerate}[label=\alph*.)]
  \item 证明：当 $\beta\to\infty$ 时，

\begin{equation}
  \langle\phi_2|e^{-\frac{\beta H}{\hbar}}|\phi_1\rangle\sim 2e^{\frac{\mathrm{i}}{2}(\phi_1-\phi_2)}\cos\left(\frac{\phi_1-\phi_2}{2}\right)e^{-\beta E_{\mathrm{ground}}}+\cdots
  \tag{2.26}
  \label{eq:2.26}
\end{equation}

  \item 注意这个表达式在 $\phi\to-\phi$ 下并不保持不变。但在 $\mathrm{Pin}^{+}(2)$ 中存在一个元素 $P$，它对应 $\phi\to-\phi$。这与我们关于 $\mathrm{Pin}^{+}(2)$ 是量子理论的一个合法对称性的论证如何相容？\footnote{\emph{答}：见 [48]。}
\end{enumerate}

\subsubsection{对整体 $SO(2)$ 对称性作规范化与 Chern-Simons 项} % 原文: 2.1.3 Gauging The Global $SO(2)$ Symmetry And Chern-Simons Terms

% [Tutor Note: 术语取舍。「gauging a symmetry」在中文文献里有「规范化」「局域化」「引规范」「把对称性变成规范对称性」等说法，本译本统一作「规范化」，并在首次出现处括注英文；但须与「normalization（归一化）」区分——同一页上（原 p.19 起）作者反复讲 normalizations 指的是耦合常数与曲率的归一化，两者不可混同。另外，这里的「gauging」是两步：先局域化并耦合规范场，再对规范场积分；只做第一步时称「与（非动力学）外源规范场耦合」，不要把它等同于「把规范场量子化」。]

当一个理论具有对称性时，人们可以实施一个称为「规范化」（gauging the symmetry）的程序。这是一个两步过程：

\begin{enumerate}
  \item 把对称性变成局域的，并耦合到一个规范场；
  \item 对规范场积分。
\end{enumerate}

并不需要在完成第 (1) 步之后接着做第 (2) 步。在这种情形下，我们说是在与\emph{非动力学的外源规范场}（nondynamical external gauge fields）耦合。为一个对称性引入非动力学的外源规范场是完全可以的。我们在量子力学课上一直这么干：把量子系统耦合到电磁场，却并不试图去量子化那个电磁场。

对数学修养更好的读者，这个两步过程可以概括为：给定一个具有对称群 $G$ 的场论，我们可以通过改变其定义域所在的配边范畴（bordism category）来使该理论成为 $G$-等变的：

\begin{enumerate}
  \item 把对称群等同于某个主丛的结构群，并把场论的几何定义域改为带联络的丛；
  \item 对主丛的同构类求和，并对联络积分。
\end{enumerate}

我们将在第 {\small\itshape ［原文标记：****］}% [原文待办标记：交叉引用的节号未填，原文 p.16 印作「in section ****」]
节中对此再多作一点解释。

在目前这个简单的例子里，我们可以把对所有 $\mathcal{B}$ 值都存在的整体 $SO(2)$ 对称性 $\phi\to\phi+\alpha$ 加以「规范化」。\footnote{于是，看看外源规范场如何告诉我们关于电荷共轭对称性的信息，就很有意思了。关于这一点的一个讨论见 D. Gaiotto、A. Kapustin、Z. Komargodski 与 N. Seiberg，\emph{Theta, Time Reversal, and Temperature}，\url{https://arxiv.org/pdf/1703.00501.pdf} 的附录 D。}

于是，在我们这个简单例子里，我们把这个平移对称性变成局域的，也就是说，我们试图把

\begin{equation}
  \phi(t)\to\phi(t)+\alpha(t)
  \tag{2.27}
  \label{eq:2.27}
\end{equation}

变成一个对称性，其中 $\alpha(t)$ 不是常数，而是时间的「任意」函数。当然，作用量 $\sim\int\dot\phi^{2}$ 在这样的变换下并不不变。我们引入一个规范场 $A_t^{(e)}dt$，其中上标 $e$ 表示「外源的」（external），提醒我们这\underline{不是}电磁场的规范场。（那个场已经产出了我们的 theta 项。）相反，它是我们系统中的一个新场：转动对称性的规范场。

规范化后的作用量是

\begin{equation}
  S=\int\frac{1}{2}I(\dot\phi+A_t^{(e)})^{2}dt+\oint\mathcal{B}(\dot\phi+A_t^{(e)})dt
  \tag{2.28}
  \label{eq:2.28}
\end{equation}

它在如下规范变换下是规范不变的：

\begin{gather}
  \phi\to\phi+\alpha(t)\notag\\
  A_t^{(e)}\to A_t^{(e)}-\partial_t\alpha(t)
  \tag{2.29}
  \label{eq:2.29}
\end{gather}

或者，更好的写法是

\begin{gather}
  e^{\mathrm{i}\phi(t)}\to e^{\mathrm{i}\phi(t)}e^{\mathrm{i}\alpha(t)}\notag\\
  A_t^{(e)}\to A_t^{(e)}+\mathrm{i}e^{-\mathrm{i}\alpha(t)}\partial_te^{\mathrm{i}\alpha(t)}
  \tag{2.30}
  \label{eq:2.30}
\end{gather}

这样更好，是因为在拓扑非平凡的时空上工作时——例如在欧氏时间圆上——应当满足单值性的是 $e^{\mathrm{i}\alpha(t)}$。

在没有外源规范场 $A_t^{(e)}$ 时，我们发现 $\mathcal{B}$ 是模一个整数周期的。往作用量里加一个 Chern-Simons 项，我们就能把某种周期性还给 $\mathcal{B}$。在欧氏空间中，新的作用量是：

\begin{equation}
  e^{-S}=e^{-\int\frac{1}{2}I(\dot\phi+A_t^{(e)})^{2}dt-\mathrm{i}\oint\mathcal{B}(\dot\phi+A_t^{(e)})dt}\,e^{\mathrm{i}k\int A_t^{(e)}dt}
  \tag{2.31}
  \label{eq:2.31}
\end{equation}

我们手工地把不变性恢复了出来，使得带有 $(\mathcal{B},k)$ 的新理论等价于带有 $(\mathcal{B}+1,k+1)$ 的理论。但现在我们必须操心新加这一项的规范不变性。这里有一个关键点：

\begin{center}
  \fbox{\parbox{0.92\textwidth}{\itshape 作用量本身不必保持不变。要有一个良定义的路径积分，所需的全部只是：指数化后的作用量必须不变。}}
\end{center}

我们后面还会多次援引这一观察：在量子化作用量中其它 Chern-Simons 项的系数时，以及在讨论二维 CFT 中的 WZ 项时。我们将称之为\emph{多值作用量原理}（multi-valued action principle）：作用量 $S$ 可以是多值的，只要 $e^{\mathrm{i}S}$ 是单值的。

因此，我们来研究

\begin{equation}
  e^{\mathrm{i}k\oint A_t^{(e)}dt}
  \tag{2.32}
  \label{eq:2.32}
\end{equation}

的规范不变性。

现在我们必须决定自己在哪个 1-流形上工作。

在区间 $[t_1,t_2]$ 上，表达式

\begin{equation}
  \exp\left[\mathrm{i}\int_{t_1}^{t_2}A_t^{(e)}(t')dt'\right]
  \tag{2.33}
  \label{eq:2.33}
\end{equation}

不是规范不变的。更确切地说：

\begin{equation}
  \exp\left[\mathrm{i}\int_{t_1}^{t_2}A_t^{(e)}(t')dt'\right]\to e^{-\mathrm{i}\alpha(t_2)}\exp\left[\mathrm{i}\int_{t_1}^{t_2}A_t^{(e)}(t')dt'\right]e^{\mathrm{i}\alpha(t_1)}
  \tag{2.34}
  \label{eq:2.34}
\end{equation}

因此，在实直线上，只要选取边界条件使得 $\alpha(t)$ 当 $t\to\pm\infty$ 时趋于零，$\exp\left[\mathrm{i}\int_{\mathbb{R}}A_t^{(e)}(t')dt'\right]$ 就是规范不变的。

类似地，在圆上，如果 $\alpha(t)$ 是单值的，那么

\begin{equation}
  \exp\left[\mathrm{i}\oint_{S^{1}}A_t^{(e)}(t')dt'\right]
  \tag{2.35}
  \label{eq:2.35}
\end{equation}

是规范不变的。然而对圆来说，真正的规范参数是 $g(t)=e^{\mathrm{i}\alpha(t)}$（对 $U(1)$ 规范群而言）。因此，$\alpha(t)=2\pi wt/\beta$（$w\in\mathbb{Z}$）是一个合法的规范变换；但在这样的规范变换下：

\begin{equation}
  \exp\left[\mathrm{i}k\oint_{S^{1}}A_t^{(e)}(t')dt'\right]\to e^{2\pi\mathrm{i}wk}\exp\left[\mathrm{i}k\oint_{S^{1}}A_t^{(e)}(t')dt'\right]
  \tag{2.36}
  \label{eq:2.36}
\end{equation}

因此，如果我们想要让理论在圆上说得通，那么 $k$ 就应当被量子化为整数。一般地：$t\mapsto e^{\mathrm{i}\alpha(t)}$ 定义了一个映射 $S^{1}_{\mathrm{time}}\to U(1)$。但 $\pi_1(U(1),*)\cong\mathbb{Z}$ 由绕数刻画。当这个绕数非零时，我们说 $e^{\mathrm{i}\alpha(t)}$ 定义了一个\emph{大的规范变换}（large gauge transformation）。

{\itshape「Chern-Simons 项」在大的规范变换下的规范不变性蕴含耦合常数是量子化的：$k\in\mathbb{Z}$。\par}

出于与电荷共轭对称性的实现有关的理由，把定义延拓到半整数的 $k$ 是有好处的。有办法赋予半整数量子化的 Chern-Simons 项以意义，那就是把 $0+1$ 维理论看作某个良定义的 $1+1$ 维理论的边界。由 Stokes 定理我们有：

\begin{equation}
  \exp\left[\mathrm{i}k\oint_{S^{1}}A_t^{(e)}dt\right]=\exp\left[\mathrm{i}k\int_{\Sigma}F^{(e)}\right]
  \tag{2.37}
  \label{eq:2.37}
\end{equation}

其中 $F^{(e)}=dA^{(e)}$。右端即使在 $k$ 不是整数时也说得通，但现在这个表达式依赖于 $1+1$ 维时空 $\Sigma$ 的「体内」（bulk）规范场的细节。

在真实的凝聚态系统中可以观测到一个非常类似的现象：$3+1$ 维拓扑绝缘体的边界理论，由能级为半整数的 Chern-Simons 理论描述。（也就是说，是朴素规范不变性所允许的能级的一半。）

\medskip\noindent\textbf{练习}\quad\emph{谜题}

解决如下悖论：

我们首先论证了：若 $k\notin\mathbb{Z}$，则 \eqref{eq:2.37} 的左端在大的规范变换下\underline{不是}不变的。然后我们又用右端的表达式来\emph{定义}左端，而右端是明显规范不变的。这两个论断如何相容？\footnote{\emph{答}：规范变换 $e^{\mathrm{i}\alpha(t)}$ 必须延拓为连续映射 $\Sigma\to U(1)$。若 $\Sigma$ 是一个光滑流形，其边界只有 $S^{1}$——这正是我们写下方程 \eqref{eq:2.37} 时默认的假设——那么这样的映射限制到作为其边界的 $S^{1}$ 上时，总是小规范变换。}

% [Tutor Note: 这条悖论的解答完全在于「延拓」二字（见上一条脚注），而不是在于 (2.37) 的等式本身：左端之所以能为半整数 $k$ 定义，靠的是把 $1+1$ 维的 $\Sigma$ 一并纳入定义，此时规范变换的参数 $\alpha$ 必须是 $\Sigma$ 上的整体定义的单值函数，它在边界 $S^{1}$ 上的绕数必为零，于是 (2.36) 里的 $e^{2\pi\mathrm{i}wk}$ 根本不会被激发。代价是：真正的半整数 Chern-Simons 项还需要额外的数据（旋量结构），本讲义后文（第 2.5 节；另见下一片段的 \eqref{eq:2.50} 前后）才展开这一点，此处不宜把半整数 $k$ 当成「无需额外数据」的结果使用。]

\subsection{三维中的 $U(1)$ Chern-Simons 理论} % 原文: 2.2 $U(1)$ Chern-Simons Theory In 3 Dimensions

\subsubsection{$U(1)$ 规范理论的若干预备知识} % 原文: 2.2.1 Some $U(1)$ Gauge Theory Preliminaries

在整份讲义中，我们将需要非常小心地对待归一化与周期性，所以我们把 $U(1)$ 规范理论的归一化明确写出来。

我们从某个时空 $M$ 上基于阿贝尔 Lie 群 $U(1)$ 或 $\mathbb{R}$ 的规范理论出发。局部地，规范场是一个实 1-形式，记作 $A=A_\mu dx^\mu$。我们要强调：虽然 $F=dA$ 是一个整体良定义的 2-形式，但一般来说 $A$ 并不是整体良定义的 1-形式，所以 $F$ 完全可能是一个非恰当（non-exact）的闭 2-形式。

更正式地说，$A$ 是 $M$ 上某个主 $U(1)$ 丛或复线丛上联络的局部 1-形式表示。在我们的归一化下，协变导数局部地写作 $D=d+\mathrm{i}A$，并在两片覆盖的交叠 $\mathcal{U}_\alpha\cap\mathcal{U}_\beta$ 上按 $D_\alpha=g_{\alpha\beta}^{-1}D_\beta g_{\alpha\beta}$ 变化。

\textbf{Remark：}\underline{一个重要的约定。}尽管从数学家的观点看这并非最好的归一化，我们这里仍取（局部定义的 1-形式）$A$ 在规范群为 $G=U(1)$ 时是一个\underline{实} 1-形式。当 $G=U(1)^{d}$ 时，我们继续取它是 $d$ 个局部定义 1-形式组成的 $d$-重 $A^{I}$。然而，在处理一般的非阿贝尔 Lie 群时，我们取 $A$ 是一个取值于 Lie 代数 $\mathfrak{g}$ 的 1-形式。现在，由于我们处理的是紧 Lie 群，若 $G$ 是一个矩阵群，则 $\mathfrak{g}$ 将是某个\underline{反厄米}矩阵 Lie 代数的子代数。在这一约定下，协变导数为 $D_A=d+A$。

如上所述，由于 $A$ 只是局部定义的，2-形式 $F$ 虽然整体有定义，却未必恰当。我们的规范场这样归一化，使得 $F/2\pi$ 具有整数周期，也就是说，绕闭 2-圈的积分总是整数：

\begin{equation}
  \int_{\Sigma_2}F\in 2\pi\mathbb{Z}
  \tag{2.38}
  \label{eq:2.38}
\end{equation}

对任何闭的定向 2-圈 $\Sigma_2\subset M$ 成立。此外，若规范群是 $G=U(1)$，则总存在流形 $M$、线丛 $L\to M$ 以及 2-圈 $\Sigma_2\subset M$，使得整数 $\int_{\Sigma_2}\frac{F}{2\pi}$ 恰好等于一。因此这是最优的归一化。反之，若规范群是 $G=\mathbb{R}$，则 $F$ 的周期将总是零。也就是说，$F$ 总是恰当的，而 $A$ 可以取为整体良定义的 1-形式。

在上述约定下，规范变换是这样的：协变导数 $D=d+\mathrm{i}A$ 按 $D\to g^{-1}Dg$ 变换，其中

\begin{equation}
  g:M\to U(1)
  \tag{2.39}
  \label{eq:2.39}
\end{equation}

是一个规范变换。如果我们能取对数，并对一个整体良定义的函数取 $g=e^{\mathrm{i}\epsilon}$，

\begin{equation}
  \epsilon:M\to\mathbb{R}
  \tag{2.40}
  \label{eq:2.40}
\end{equation}

那么

\begin{equation}
  A\to A-d\epsilon
  \tag{2.41}
  \label{eq:2.41}
\end{equation}

我们把这样的变换称为「小的规范变换」（small gauge transformations）。若规范群是 $\mathbb{R}$，则所有规范变换都是这个形式。另一方面，若规范群是 $U(1)$，则会出现不允许整体对数的映射 $g:M\to U(1)$。一个简单的判据是：1-形式 $\omega=-\mathrm{i}g^{-1}dg$ 是否有非零的周期。如果 $g(x)$ 不能写成 $g(x)=e^{\mathrm{i}\epsilon(x)}$（其中 $\epsilon(x)$ 是整体良定义的函数），我们就说 $g(x)$ 是一个「大的规范变换」。特别地，若 $\pi_1(M,x_0)$ 非平凡，则存在大的规范变换。在这种情形下，规范变换最好被设想为平移

\begin{equation}
  A\to A+\omega
  \qquad \omega\in\Omega^{1}_{2\pi\mathbb{Z}}(M)
  \tag{2.42}
  \label{eq:2.42}
\end{equation}

其中 $\Omega^{1}_{2\pi\mathbb{Z}}(M)$ 是所有（适当可微的）周期都落在 $2\pi\mathbb{Z}$ 中的闭 1-形式构成的空间。

% [Tutor Note: (2.41) 的符号必须与约定绑在一起核对。按正文自己给的 $D=d+\mathrm{i}A$、$D\to g^{-1}Dg$、$g=e^{\mathrm{i}\epsilon}$ 直接算：$g^{-1}(\partial+\mathrm{i}A)g=\partial+g^{-1}\partial g+\mathrm{i}A$，而 $g^{-1}\partial g=\mathrm{i}\partial\epsilon$，于是得到 $A\to A+d\epsilon$，与印出的 (2.41)「$A\to A-d\epsilon$」差一个符号。把相邻几式放在一起看会更清楚：(2.30) 的 $A_t^{(e)}\to A_t^{(e)}+\mathrm{i}e^{-\mathrm{i}\alpha}\partial_te^{\mathrm{i}\alpha}$ 其实等于 $A_t^{(e)}-\partial_t\alpha$，即 $\delta A=-d\alpha$；而 (2.45) 的右端 $d(\epsilon F)$ 在 $dF=0$ 时等于 $d\epsilon\wedge F$，对应的却是 $\delta A=+d\epsilon$（若与 (2.41) 配套，本应印成 $A\wedge F-d(\epsilon F)$）。也就是说，(2.30) 与 (2.41) 自相一致，而 (2.45) 取的是相反的一支；这源于原稿在「$A$ 取实值」与「$A$ 取反厄米」两种归一化之间往返。复算 $\kappa$ 的量子化条件、或把 (2.48) 的整体符号与别的文献对照时，务必先固定一支（例如 $\delta A=-d\epsilon$、$\omega=-\mathrm{i}g^{-1}dg$），不要跨式混用。]

此外，当 $M$ 非紧或者有非空边界时，我们还需要讨论规范场与规范群的边界条件。我们会考虑不同种类的边界条件，所以这一点留到后面再说。

\subsubsection{从 $\theta$ 项到 Chern-Simons} % 原文: 2.2.2 From $\theta$-term To Chern-Simons

现在，由于 $F=dA$，我们注意到四维的拓扑密度可以写成：

\begin{equation}
  F\wedge F=d(AdA)
  \tag{2.43}
  \label{eq:2.43}
\end{equation}

如果我们考虑一个 $\partial M_4=M_3$ 的四维路径积分（$M_3$ 可以连通，也可以不连通），那么我们就会被要求把

\begin{equation}
  \int_{M_3}AdA
  \tag{2.44}
  \label{eq:2.44}
\end{equation}

作为作用量中的一项来考虑。

初看起来，\eqref{eq:2.44} 不像是一个能放进规范理论作用量里的项：这一项明显含有 $A$ 这个因子，因此显然不是规范不变的。另一方面，$F\wedge F$ 是规范不变的，所以 $AdA$ 的变化必须是 $d$-闭的。我们可以指望它是 $d$-恰当的，从而「由全导数造成的变化无关紧要」。

那么，我们来算一算：

\begin{enumerate}
  \item 在小规范变换 \eqref{eq:2.41} 下，我们有

\begin{equation}
  A\wedge F\to A\wedge F+d(\epsilon F)
  \tag{2.45}
  \label{eq:2.45}
\end{equation}

  这确实是一个全导数，正如我们所希望的那样。如果忽略边界贡献（一般来说我们\underline{不能}这么做），那么规范不变性就有保证了。无论如何，这是一个好的开始。

  \item 在大规范变换 \eqref{eq:2.42} 下，我们有

\begin{equation}
  A\wedge F\to A\wedge F+\omega\wedge F
  \tag{2.46}
  \label{eq:2.46}
\end{equation}

  现在，任何具有量子化周期的 1-形式必定是闭 1-形式，但一般来说它不必整体恰当，所以 $\omega\wedge F$ 不必整体恰当。连同边界项的问题，这一点是我们需要处理的。
\end{enumerate}

\medskip\noindent\underline{$\kappa$ 的量子化}% 原文: Quantization Of $\kappa$（原以斜体下划线标示的无编号小标题，独占一行）
\par\vspace{2pt}

我们考虑在某个定向 3-流形 $M_3$ 上的 $U(1)$ 规范理论作用量中加入可能的 Chern-Simons 项，并把耦合常数归一化为：

\begin{equation}
  \exp\left[\mathrm{i}\frac{\kappa}{2\pi}\int_{M_3}A\wedge F\right]
  \tag{2.47}
  \label{eq:2.47}
\end{equation}

暂时我们取 $M_3$ 没有边界。即便在这种情况下，正如我们刚刚讨论过的，作用量并不是规范不变的。在路径积分中，我们并不对 $\mathcal{A}$——所有规范场（即固定主丛上的所有联络）构成的空间——积分，而是对 $\mathcal{A}/\mathcal{C}G$——所有规范等价类的空间——积分。这里 $\mathcal{G}$ 是规范变换构成的群。通过 Fadeev-Popov 或 BRST 型的程序，我们可以把 $\mathcal{A}$ 上的自然测度 $[dA]$ 推下来，成为 $\mathcal{A}/\mathcal{G}$ 上的一个测度；但这样一来，当我们用 $e^{\mathrm{i}S}$ 去权重这个测度时，那个因子必须是空间 $\mathcal{A}/\mathcal{G}$ 上的良定义且单值的函数。

当你看到作用量里出现一个显式的规范场时，你脑子里的所有警报都应当响起来：这个理论多半说不通。Chern-Simons 之所以能躲过这个问题，靠的就是上面提到的多值作用量原理。这个原理之所以适用，是因为该理论的路径积分是对 $\mathcal{A}/\mathcal{G}$（规范场的规范等价类空间）的积分，而我们所需要的一切，只是这个空间上的一个（形式上）良定义的测度 $[dA]e^{\mathrm{i}S}$。

让我们回到 \eqref{eq:2.46}。那么，正如我们刚刚指出的，Chern-Simons 作用量 $\int_{M_3}AdA$ \underline{不是}规范不变的。在 \eqref{eq:2.46} 下，指数化作用量的变化是一个乘性因子：

\begin{equation}
  e^{2\pi\mathrm{i}\frac{\kappa}{4\pi^{2}}\int_{M_3}A\wedge dA}\to e^{2\pi\mathrm{i}\frac{\kappa}{4\pi^{2}}\int_{M_3}A\wedge dA}\cdot e^{2\pi\mathrm{i}\kappa\int_{M_3}\frac{\omega}{2\pi}\wedge\frac{F}{2\pi}}
  \tag{2.48}
  \label{eq:2.48}
\end{equation}

于是，如果因子

\begin{equation}
  e^{2\pi\mathrm{i}\kappa\int_{M_3}\frac{\omega}{2\pi}\wedge\frac{F}{2\pi}}
  \tag{2.49}
  \label{eq:2.49}
\end{equation}

对所有规范变换 $\omega\in\Omega^{1}_{2\pi\mathbb{Z}}(M_3)$ 以及所有规范场都总是等于一，那么我们就能确保路径积分测度的规范不变性。但是，再一次强调，这恰恰就是我们为了得到良定义的路径积分、从而得到良定义的理论所需要的全部。（至少当 $M_3$ 没有边界时是这样。）

按照假设，我们取 $\omega\in\Omega^{1}_{2\pi\mathbb{Z}}(M_3)$ 为任意的、周期在 $2\pi\mathbb{Z}$ 中的闭 1-形式。现在，拓扑学中的一个定理 [12, 51] 向我们保证：若规范群是 $U(1)$，则形式 $F/2\pi$ 具有任意的整数周期。因此，路径积分测度是良定义的，当且仅当 $\kappa\in\mathbb{Z}$。

% [Tutor Note: (2.48)–(2.49) 里的 $2\pi$ 账目值得逐项核一遍，这是最容易漏掉因子的地方：由 (2.47) 的 $\frac{\kappa}{2\pi}\int A\wedge F$ 及 $F=dA$ 可得 $2\pi\mathrm{i}\frac{\kappa}{4\pi^{2}}=\frac{\mathrm{i}\kappa}{2\pi}$，故 (2.48) 左端的指数就是 (2.47) 的指数；而 (2.46) 的平移 $A\wedge F\to A\wedge F+\omega\wedge F$ 给出的相位增量正是 $\frac{\mathrm{i}\kappa}{2\pi}\int\omega\wedge F=2\pi\mathrm{i}\kappa\int\frac{\omega}{2\pi}\wedge\frac{F}{2\pi}$，与 (2.48) 右端的乘性因子、(2.49) 逐字吻合。另注：本段正文把规范等价类空间印作 $\mathcal{A}/\mathcal{C}G$（原 p.21），多出的 $\mathcal{C}$ 应是排版笔误，同段下文又写回 $\mathcal{A}/\mathcal{G}$；译文照原文转写未改，读者应以 $\mathcal{A}/\mathcal{G}$ 为准。]

\medskip\noindent\textbf{Remarks（评注）：}
